\documentclass[11pt]{article}
\usepackage[margin=1in]{geometry}
\usepackage{amsmath,amssymb,booktabs,array,longtable}
\usepackage{graphicx}
\usepackage{hyperref}
\usepackage{enumitem}
\usepackage{titlesec}
\usepackage{float}
\begin{document}

\begin{center}
\Large\textbf{Shiv Nadar University Chennai}\\
\end{center}


\begin{center}
\textbf{CS3807 – Deep Learning Laboratory}
\vspace{2mm}

\textbf{Experiment 3}
\vspace{2mm}

\textbf{Implementation of Convolutional Neural Networks (CNNs) for Image Classification}
\end{center}


\renewcommand{\arraystretch}{1.4}

\begin{tabular}{|p{3.5cm}|p{8cm}|p{2cm}|}
\hline
Degree \& Branch & B.Tech Artificial Intelligence \& Data Science & Semester V\\ \hline
Subject Code \& Name & CS3807 -- Deep Learning Laboratory & AY: 2026--27\\ \hline
\end{tabular}
\vspace{3mm}

\section{Objective}
To understand the working principle of Convolutional Neural Networks by implementing convolution, pooling, feature map visualization, and image classification using TensorFlow/Keras.

\section{Learning Outcomes}
Students will be able to:
\begin{enumerate}[leftmargin=*]
    \item Explain the convolution operation.
    \item Compute output dimensions using stride and padding.
    \item Visualize feature maps.
    \item Implement max pooling and average pooling.
    \item Calculate CNN parameters.
    \item Build a CNN for image classification.
    \item Evaluate CNN performance.
\end{enumerate}

\section{Background Theory}
A Convolutional Neural Network consists of:
\begin{itemize}
    \item Convolution Layer
    \item Activation Layer
    \item Pooling Layer
    \item Fully Connected Layer
\end{itemize}

The convolution operation is:
\[
Y(i,j) = \sum_{m}\sum_{n} X(i+m,j+n) K(m,n)
\]
where:
\begin{itemize}[label={}]
    \item $X$: Input Image
    \item $K$: Kernel (Filter)
    \item $Y$: Feature Map
\end{itemize}

The output feature map size is:
\[
\frac{N - F + 2P}{S} + 1
\]
where:
\begin{itemize}[label={}]
    \item $N$: Input Size
    \item $F$: Filter Size
    \item $P$: Padding
    \item $S$: Stride
\end{itemize}

\subsection{Common Convolution Kernels}
A kernel (or filter) is a small matrix that slides over an input image to extract important visual features such as edges, textures, corners, and smooth regions. During convolution, the kernel is multiplied element-wise with the corresponding image pixels, and the results are summed to produce a single value in the output feature map. Different kernels highlight different image characteristics.

\subsubsection*{1. Identity Kernel}
The identity kernel preserves the original image without modifying its pixel values.
\[
\begin{bmatrix}
0 & 0 & 0 \\
0 & 1 & 0 \\
0 & 0 & 0
\end{bmatrix}
\]
\textbf{Purpose:}
\begin{itemize}
    \item Preserves the original image.
    \item Used for understanding the convolution operation.
\end{itemize}

\subsubsection*{2. Sobel-X Kernel (Vertical Edge Detection)}
The Sobel-X kernel detects vertical edges by measuring intensity changes along the horizontal direction.
\[
\begin{bmatrix}
-1 & 0 & 1 \\
-2 & 0 & 2 \\
-1 & 0 & 1
\end{bmatrix}
\]
\textbf{Applications:}
\begin{itemize}
    \item Object detection
    \item Lane detection
    \item Face recognition
\end{itemize}

\subsubsection*{3. Sobel-Y Kernel (Horizontal Edge Detection)}
The Sobel-Y kernel detects horizontal edges by measuring intensity changes along the vertical direction.
\[
\begin{bmatrix}
-1 & -2 & -1 \\
0 & 0 & 0 \\
1 & 2 & 1
\end{bmatrix}
\]
\textbf{Applications:}
\begin{itemize}
    \item Text detection
    \item Boundary extraction
    \item Medical image analysis
\end{itemize}

\subsubsection*{4. Laplacian Kernel}
The Laplacian kernel detects edges in all directions using the second-order derivative of image intensity.
\[
\begin{bmatrix}
0 & -1 & 0 \\
-1 & 4 & -1 \\
0 & -1 & 0
\end{bmatrix}
\]
\textbf{Applications:}
\begin{itemize}
    \item Edge enhancement
    \item Satellite image analysis
    \item Medical imaging
\end{itemize}

\subsubsection*{5. Sharpening Kernel}
This kernel enhances fine details by emphasizing high-frequency components.
\[
\begin{bmatrix}
0 & -1 & 0 \\
-1 & 5 & -1 \\
0 & -1 & 0
\end{bmatrix}
\]
\textbf{Applications:}
\begin{itemize}
    \item Image enhancement
    \item Optical Character Recognition (OCR)
    \item Face recognition
\end{itemize}

\subsubsection*{6. Box Blur Kernel}
The Box Blur kernel smooths an image by replacing each pixel with the average of its neighboring pixels.
\[
\frac{1}{9}
\begin{bmatrix}
1 & 1 & 1 \\
1 & 1 & 1 \\
1 & 1 & 1
\end{bmatrix}
\]
\textbf{Applications:}
\begin{itemize}
    \item Noise removal
    \item Image smoothing
    \item Image preprocessing
\end{itemize}

\subsubsection*{7. Gaussian Blur Kernel}
Gaussian blur smooths the image while preserving the overall structure better than a simple average filter.
\[
\frac{1}{16}
\begin{bmatrix}
1 & 2 & 1 \\
2 & 4 & 2 \\
1 & 2 & 1
\end{bmatrix}
\]
\textbf{Applications:}
\begin{itemize}
    \item Noise reduction
    \item Computer vision preprocessing
    \item Feature extraction
\end{itemize}

\subsubsection*{8. Emboss Kernel}
The Emboss kernel creates a three-dimensional appearance by emphasizing directional changes in intensity.
\[
\begin{bmatrix}
-2 & -1 & 0 \\
-1 & 1 & 1 \\
0 & 1 & 2
\end{bmatrix}
\]
\textbf{Applications:}
\begin{itemize}
    \item Texture analysis
    \item Artistic image effects
\end{itemize}

\subsubsection*{9. Outline Detection Kernel}
This kernel highlights object boundaries by emphasizing regions with rapid intensity changes.
\[
\begin{bmatrix}
-1 & -1 & -1 \\
-1 & 8 & -1 \\
-1 & -1 & -1
\end{bmatrix}
\]
\textbf{Applications:}
\begin{itemize}
    \item Image segmentation
    \item Boundary detection
    \item Shape extraction
\end{itemize}

\subsubsection*{10. Motion Blur Kernel}
This kernel simulates motion blur along the diagonal direction.
\[
\frac{1}{5}
\begin{bmatrix}
1 & 0 & 0 & 0 & 0 \\
0 & 1 & 0 & 0 & 0 \\
0 & 0 & 1 & 0 & 0 \\
0 & 0 & 0 & 1 & 0 \\
0 & 0 & 0 & 0 & 1
\end{bmatrix}
\]
\textbf{Applications:}
\begin{itemize}
    \item Motion simulation
    \item Image restoration
    \item Video processing
\end{itemize}

\section{Summary of Common Kernels}
\begin{table}[h!]
\centering
\begin{tabular}{|l|l|l|l|}
\hline
\textbf{Kernel} & \textbf{Size} & \textbf{Purpose} & \textbf{Applications} \\ \hline
Identity & $3\times3$ & Preserve original image & Baseline comparison \\ \hline
Sobel-X & $3\times3$ & Detect vertical edges & Object detection \\ \hline
Sobel-Y & $3\times3$ & Detect horizontal edges & Boundary detection \\ \hline
Laplacian & $3\times3$ & Detect edges in all directions & Medical imaging \\ \hline
Sharpen & $3\times3$ & Enhance details & Image enhancement \\ \hline
Box Blur & $3\times3$ & Smooth image & Noise removal \\ \hline
Gaussian Blur & $3\times3$ & Reduce noise & Computer vision \\ \hline
Emboss & $3\times3$ & 3D effect & Artistic filtering \\ \hline
Outline & $3\times3$ & Boundary extraction & Segmentation \\ \hline
Motion Blur & $5\times5$ & Simulate motion & Video processing \\ \hline
\end{tabular}
\end{table}

\section{Numerical Example 1: Convolution}
Consider the following input image:
\[
X = \begin{bmatrix} 1 & 2 & 3 \\ 4 & 5 & 6 \\ 7 & 8 & 9 \end{bmatrix}
\]
Kernel:
\[
K = \begin{bmatrix} 1 & 0 \\ 0 & 1 \end{bmatrix}
\]
Output:
\begin{itemize}[label={}]
    \item \textbf{First position:} $1(1) + 2(0) + 4(0) + 5(1) = 6$
    \item \textbf{Second position:} $2(1) + 3(0) + 5(0) + 6(1) = 8$
    \item \textbf{Third position:} $4(1) + 5(0) + 7(0) + 8(1) = 12$
    \item \textbf{Fourth position:} $5(1) + 6(0) + 8(0) + 9(1) = 14$
\end{itemize}
Feature Map:
\[
\begin{bmatrix} 6 & 8 \\ 12 & 14 \end{bmatrix}
\]

\section{Numerical Example 2: Max Pooling}
Feature map:
\[
\begin{bmatrix}
1 & 5 & 2 & 3 \\
7 & 8 & 1 & 0 \\
4 & 6 & 9 & 5 \\
2 & 3 & 1 & 8
\end{bmatrix}
\]
Using a $2 \times 2$ max pooling filter:
\begin{itemize}[label={}]
    \item \textbf{Window 1:} $\begin{bmatrix} 1 & 5 \\ 7 & 8 \end{bmatrix} \rightarrow 8$
    \item \textbf{Window 2:} $\begin{bmatrix} 2 & 3 \\ 1 & 0 \end{bmatrix} \rightarrow 3$
    \item \textbf{Window 3:} $\begin{bmatrix} 4 & 6 \\ 2 & 3 \end{bmatrix} \rightarrow 6$
    \item \textbf{Window 4:} $\begin{bmatrix} 9 & 5 \\ 1 & 8 \end{bmatrix} \rightarrow 9$
\end{itemize}
Output:
\[
\begin{bmatrix} 8 & 3 \\ 6 & 9 \end{bmatrix}
\]

\section{Numerical Example 3: Parameter Calculation}
Suppose:
\begin{itemize}
    \item \textbf{Input Image:} $32 \times 32 \times 3$
    \item \textbf{Convolution Layer:} Filters $= 16$, Kernel $= 3 \times 3$
\end{itemize}
Parameters:
\[
(3 \times 3 \times 3 + 1) \times 16 = (27 + 1) \times 16 = 448
\]
Therefore,
\[
\text{Total trainable parameters} = 448.
\]

\section{Dataset}
\textbf{Dataset: CIFAR-10}
\begin{itemize}
    \item Training Images: 50,000
    \item Testing Images: 10,000
    \item Classes: 10
    \item Image Size: $32 \times 32 \times 3$
\end{itemize}

\section{Experimental Procedure}
\begin{description}
    \item[Task 1] Load CIFAR-10.
    \begin{itemize}
        \item Display ten sample images.
        \item Print dataset dimensions.
        \item Plot class distribution.
    \end{itemize}

    \item[Task 2] Implement a convolution layer. Compare:
    \begin{itemize}
        \item Kernel $= 3 \times 3$
        \item Kernel $= 5 \times 5$
        \item Kernel $= 7 \times 7$
    \end{itemize}
    Record feature map sizes.

    \item[Task 3] Study hyperparameters. Compare:
    \begin{itemize}
        \item Stride $= 1$
        \item Stride $= 2$
        \item Padding $=$ Same
        \item Padding $=$ Valid
    \end{itemize}
    Compute output dimensions.

    \item[Task 4] Visualize feature maps after the first convolution layer. Display at least eight feature maps.

    \item[Task 5] Compare:
    \begin{itemize}
        \item Max Pooling
        \item Average Pooling
    \end{itemize}
    Compare output size and accuracy.

    \item[Task 6] Construct the following CNN:
    \[
    \text{Input} \rightarrow \text{Conv} \rightarrow \text{ReLU} \rightarrow \text{MaxPool} \rightarrow \text{Conv} \rightarrow \text{ReLU} \rightarrow \text{MaxPool} \rightarrow \text{Flatten} \rightarrow \text{Dense} \rightarrow \text{Softmax}
    \]
    Train for:
    \begin{itemize}
        \item Optimizer: Adam
        \item Epochs: 20
        \item Batch Size: 32
    \end{itemize}

    \item[Task 7] Evaluate:
    \begin{itemize}
        \item Accuracy
        \item Precision
        \item Recall
        \item F1-score
        \item Confusion Matrix
        \item Classification Report
    \end{itemize}
\end{description}

\section{Mandatory Plots}
Generate:
\begin{enumerate}
    \item Sample Images
    \item Class Distribution
    \item Training Accuracy
    \item Validation Accuracy
    \item Training Loss
    \item Validation Loss
    \item Feature Maps
    \item Confusion Matrix
\end{enumerate}
Write a 2--3 line inference for each plot.

\section{Additional Exercises}
\begin{enumerate}
    \item Calculate the output size for a $64 \times 64$ image using a $5 \times 5$ kernel with stride 2 and padding 2.
    \item Compute the trainable parameters of a convolution layer having 64 filters of size $3 \times 3$ and an RGB input.
    \item Compare ReLU and Sigmoid activation functions.
    \item Replace Max Pooling with Average Pooling and compare the results.
    \item Increase the number of convolution filters from 16 to 64 and analyze the effect on accuracy and computation time.
\end{enumerate}

\section{Results}
Record:
\begin{itemize}
    \item Training Accuracy
    \item Testing Accuracy
    \item Precision
    \item Recall
    \item F1-score
    \item Number of Parameters
\end{itemize}

\section{Discussion}
Answer the following:
\begin{enumerate}
    \item Why is convolution preferred over fully connected layers for images?
    \item How does stride affect the feature map size?
    \item What is the role of padding?
    \item Why is pooling used?
    \item How do feature maps represent image characteristics?
    \item Why do CNNs require fewer parameters than MLPs?
\end{enumerate}

\section{References}
\begin{enumerate}
    \item Goodfellow et al., \textit{Deep Learning}.
    \item Bishop, \textit{Pattern Recognition and Machine Learning}.
    \item Haykin, \textit{Neural Networks and Learning Machines}.
    \item TensorFlow Documentation.
    \item CIFAR-10 Dataset Documentation.
\end{enumerate}

\end{document}